% Supplementary Information of the article: deviations from the pre-registration and
% the order in which things were seen (PREREG amendments).

The hypotheses, metrics, blocks, conformal procedure and decision rules were registered (version~1.0) in the project repository before any hindcast, topography-accuracy or projection metric was computed; the repository history records the order. Dated amendments followed, each stating its reason and what had been seen. Amendments 1--12 preceded every main-run output. Amendments 13--15 came after the elevation model and training-block outputs existed but before any hindcast or test-block evaluation. Amendment~16 (generality) and Amendment~17 (ATL24) each preceded their own metrics. The deviations, in the order of the analysis:

\begin{enumerate}
\item \emph{Gauge levels before 1992.} No pre-1992 gauge series was obtained, so by the registered fallback the rising limb E1 runs from the start of altimetry to its maximum, and \metric{levels.n_scenes_by_epoch.pre_altimetry} Landsat scenes from before altimetry have no level and were not used.
\item \emph{G-REALM.} The registered cross-check series, in its NASA successor, gives level variations relative to a reference pass rather than heights on a datum\citep{gwm_atbd2025}; Hydroweb served as the cross-check.
\item \emph{Landsat only} (Amendment~12). Sentinel-2 was dropped from the main level-binned counts to stay within the Earth Engine quota; it entered the per-scene stage and one robustness variant.
\item \emph{Edge-guided threshold sample} (Amendment~3). The per-acquisition histogram is sampled inside the active zone instead of in a buffer around scene edges, at far lower cost.
\item \emph{Wind-setup screening.} Not a deviation but an outcome that limits the design: the registered rule selected no screening, and its statistic sits at the \metric{params.stage2_res_m}~m grid resolution.
\item \emph{Deep-water rule} (Amendment~13). Added after dark, clear water was found misread as land on the open shelf while the static-contour rule was checked, before any hindcast or emergence metric. The hindcast without it is red-team test R1.
\item \emph{States outside the observed cover} (Amendment~14). Widened after the first production run stopped at its built-in separation check, before any hindcast metric; the change can only make more cells unknown.
\item \emph{Validation population and block size} (Amendment~15). After training blocks only had been seen: the accuracy test uses bare emerged seabed, and, because the residual semivariogram reached no sill, the block size was set by the rule for that case and capped at \metric{h2.block_size_km}~km. The Ural delta sector then has no test block, and two sectors have one.
\item \emph{Defect in the hindcast block table.} Found after the first H1 results: undefined distances had been encoded as the smallest representable integer. It was fixed and tested, the hindcast rerun and every table compared; both runs' block tests are reported (section~\SIaudit), and no significance changed.
\item \emph{GLO-30 and GEDI} were run after the ICESat-2 decision on H2. They are registered checks, reported beside the decision; the decision is defined on ICESat-2 and was not re-made. GLO-30's acquisition window is December 2010 to January 2015, slightly wider than the registered 2011--2015.
\item \emph{Robustness checks.} The verdict rule was finalised after the tables of two variants had been seen (its magnitude criterion can only make verdicts stricter). Several variants ran on a sample of test blocks and were judged on the elevation model, not on the audit, and several registered variants were not run (section~\SIrobust). No robustness check covers H3.
\item \emph{Generality} (Amendment~16). Lake Urmia was not run within the month's Earth Engine budget, and the pre-registered substitute (the 3DEP elevation model, mostly 2016 lidar\citep{quantumspatial2017_gsl}) replaced the survey elevation model as the positive control at Great Salt Lake.
\item \emph{ATL24} (Amendment~17). Registered as the independent check below the 2026 waterline; it holds no data over the Caspian, so the check could not be made.
\item \emph{Near-term error model.} The registered \texttt{linear\_hold} result is kept as the registered H3 result. After its degenerate slopes were found, the registered sensitivity forms were run as registered; no new model form was added. Which result the article headlines was decided after the results were seen (framing: near-term emergence is limited by bathymetry), and the registered verdict is reported as registered.
\item \emph{Exposure layers.} Protected areas and Ramsar sites (registered for the exposure report, outside the decision rule) were added from the World Database on Protected and Conserved Areas after the H3 results, as exposure only, on the H3 surfaces already computed; no hypothesis or test changed. Caspian seal habitat needs a data request and was not used.
\item \emph{Display choices made after a result was seen.} The level shown for Great Salt Lake (figure~\ref{M-fig:erl_gsl} of the article) was chosen by a second rule after the first picked an atypical level; the rule used is in the repository's display code. Display choices do not enter any test.
\item \emph{Additional analyses (pre-registration Amendment 18)} (Amendment~18, registered before any of its analyses was computed). The article's headline moved to the four checked sectors, whose observations agree with ICESat-2; the near-term numbers moved from the registered primary (each grid as distributed, referenced to its own 2026 shoreline) to the grids shifted by their hindcast constants and to the registered common-baseline sensitivity, with the registered primary kept as an illustration of datum risk; the registered split-conformal intervals of emerged area are no longer presented as prediction intervals. No registered verdict was replaced, and every Amendment 18 analysis is labelled exploratory (section~\SIrevision). Three requests of the reviewers were not run, because they need Earth Engine processing beyond the study's computing budget: a split of the record by wind regime, Sentinel-2 in the lowest bins, and levelled pre-1992 Landsat scenes. The registered Great Salt Lake control could not be retrieved by script and was not downloaded by hand.
\item \emph{Analyses made in journal review (pre-registration Amendment 20)} (registered before any of its analyses was computed; section~S11). No registered verdict was replaced, and the amendment's reporting rule was not triggered: no verdict in a checked sector changed without the deltaic sectors, and the largest move of a failing grid's sector-median displacement on the Kizlyar--Dagestan and interfluve coasts was \metric{rev3.marine.four.westinterfluve.d_move_over_q_max} of the coast's noise floor (\metric{rev3.marine.four.coasts3.d_move_over_q_max} with the northeastern coast). Five departures from the amendment's text are recorded. The display rule of its item~1 named the Caspian polygon of HydroLAKES as rasterised in the Global Lakes and Wetlands Database; the polygon itself is used instead, because the \metric{rev3.band.raster.cell_arcsec} arc-second raster is too coarse on the delta fronts (the area not interpreted is \metric{rev3.band.not_seabed_km2}~\kmsq{} with the polygon and \metric{rev3.band.raster.not_seabed_km2}~\kmsq{} with the raster). The rule was written for the study-area and emergence maps of the article; it is also applied to the article's accuracy map and to Supplementary figure~\ref{fig:dem}, and the key says why the ground is not interpreted instead of naming it delta plain and floodplain, because part of it may be seabed that emerged before 2000. Item~5 specified buffers of \metric{rev3.indep.buffer_km} and \metric{rev3.indep.buffer_far_km}~km; no test point lies as far as the second from a training or calibration block, so only the first could be evaluated. Item~8 compared the wind at the acquisition hours with all hours of the same months and years; because the satellite passes at one time of day, a second reference of the same hours of the day was added, and both are reported. Two comparisons were added after the results were seen and are descriptive: the largest position error of a lidar stratum against the grids' smallest displacement on the same coast, and the largest mean setup of a level bin against the level-error budget. A manually interpreted sample of shorelines and a split of the record by wind regime need new image processing and were not run.
\end{enumerate}
